\documentclass[11pt]{report}
%               ^ Aquí puedes cambiar el tamaño de la fuente
% Imports
\usepackage[utf8]{inputenc}
\usepackage[francais]{babel}
\usepackage[T1]{fontenc}
\usepackage[margin=2cm, top=2cm, includefoot]{geometry}
\usepackage{graphicx}
\usepackage{float}
\usepackage{xcolor}
\usepackage{listings}
\usepackage{realboxes}
\usepackage{framed}
\usepackage{titlesec}
\usepackage{caption}
\usepackage{tabto}
\usepackage{url}
\usepackage{csquotes}
\usepackage{chngcntr} \counterwithout{figure}{chapter}
\usepackage[hidelinks]{hyperref}
\usepackage{parskip}
\usepackage[most]{tcolorbox}
\usepackage{fancyhdr}
\usepackage{tabularx}
\usepackage{sectsty}
\usepackage{afterpage}


\def\logouniversite{figs/upmc.png}%LOGO

% ----------------------------------------------------

% Poner "índice" en vez de "index"
\addto\captionsspanish{\renewcommand{\contentsname}{Índice}}


\setlength{\headheight}{40pt}
\pagestyle{fancy}
\fancyhf{}

% Añade una cabecera con el logo de la UPCT
\lhead{
    \includegraphics[height=40pt]{figs/Logo_Esirem.png}
    % Aquí puedes añadir contenido a tu cabecera
} 
% Puedes añadir contenido en la parte derecha de la cabecera de forma análoga con \rhead{}

% Linea negra bajo la cabecera
\renewcommand{\headrulewidth}{3pt}
\renewcommand{\headrule}{\hbox to\headwidth{\color{black}\leaders\hrule height \headrulewidth\hfill}}

% Mejor formateo de los capítulos
\titleformat{\chapter}[display]
  {\normalfont\centering\bfseries}{\vspace{-100pt}}{0pt}{\Huge}

% Mejor formateo de secciones 
\titleformat{\section}
    {\normalfont\scshape}{1em}{\thesection}{\Large}

% Función para crear hoja en blanco
\newcommand\blankpage{%
    \null
    \thispagestyle{empty}%
    \addtocounter{page}{-1}%
    \newpage}

\newcommand{\titulo}{\color{blue}Modèle épidémie SIR en épidémiologie : présentation et illustration}

\newcommand{\autor}{TURPIN Yohann, GEHU Lauric, FOSSET Jorel.}
\newcommand{\fecha}{ESIREM Dijon}


% ----------------------------------------------------

\begin{document}
    % Añadir contador de páginas en el pie de página
    \cfoot{\thepage}
    
    % Crear la portada
    \begin{titlepage}
        \reversemarginpar\marginpar{%
            \vspace{-86pt}
            \hspace*{25pt}
            \includegraphics[height=3cm]{figs/UBFC.png}%
        }
        
        {\centering
            \vspace{3cm} % edita esta distancia para subir/bajar el logo
            \hspace{5cm} \includegraphics[width=.4\textwidth]{figs/Logo_Esirem.png} %comenta esta línea si no quieres el logo en la portada
            \par\vspace{5cm}
            {\scshape\huge\textbf{\hspace{2cm}\titulo}} \par\vspace{6cm}
            
            {\LARGE
            \begin{flushright}
                    \item \autor 
                    \item \textit{\fecha}
            \end{flushright}
            }
        }
    \end{titlepage}

    % Hoja el blanco tras la portada
    \afterpage{\blankpage}
    
\begingroup
    \tableofcontents
\endgroup    
    % Inicio del texto
    \newpage
    \chapter{Le modèle SIR}
        \section{Histoire et amélioration}

	Le modèle SIR est un modèle mathématique simple qui permet de modéliser la propagation d'une maladie infectieuse dans une population. Le modèle divise la population en trois groupes : les individus susceptibles d'être infectés (S), les individus infectés (I) et les individus guéris ou retirés (R). Les individus passent d'un groupe à l'autre en fonction de leur état de santé. Le modèle est décrit par les équations différentielles suivantes :
$$\frac{dS}{dt} = -\frac{\beta SI}{N}$$

$$\frac{dI}{dt} = \frac{\beta SI}{N} - \gamma I$$

$$\frac{dR}{dt} = \gamma I$$
où S, I et R représentent respectivement le nombre de personnes dans chaque groupe, N est la taille totale de la population, $\beta$ est le taux de contact et $\gamma$ est le taux de guérison. Le modèle de Reed et Frost est un autre modèle mathématique simple qui permet de modéliser la propagation d'une maladie infectieuse dans une population. Dans ce modèle, les individus sont divisés en trois groupes : les individus susceptibles d'être infectés (S), les individus infectés (I) et les individus retirés (R). Le modèle suppose que chaque individu infecté a une probabilité fixe de transmettre la maladie à un individu susceptible chaque semaine. Le nombre de reproduction de base R0 est défini comme le nombre moyen de susceptibles qu'un individu infecté infecte au début de l'épidémie, lorsque presque toute la population est susceptible.

Le nombre de reproduction de base, R0, est défini comme le nombre attendu d'infections générées par un individu infectieux dans une grande population susceptible. Pour le modèle épidémique SIR standard que nous avons décrit précédemment, R0 est défini comme $\lambda i$ où i est la durée moyenne de la période infectieuse et $\lambda$ est l'intensité du processus de Poisson homogène qui modélise les contacts entre les individus infectés et les individus susceptibles.


            \newpage Enseñando la cabecera
		
            
            
 

\setlength{\abovedisplayskip}{6pt}
\setlength{\belowdisplayskip}{6pt}




%%% PUBLICATIONS QUE VOUS SOUHAITEZ FAIRE APPARAITRE SANS LES CITER

%Int Journal
\nocite{HKC16}
\nocite{THV14}
\nocite{HKCMST12}
%Int Conf
\nocite{FRT16a}
\nocite{FRT16}
\nocite{FRT16f}
\nocite{HdT15a}
\nocite{HdT15}
\nocite{CDR12}
\nocite{HKC12}
\nocite{CAT11a}
\nocite{CAT11b}
\nocite{wise2011}
\nocite{AAC08a}
\nocite{TDL07}
\nocite{TD07}
\nocite{DT07}
\nocite{TDL07a}
\nocite{TD07a}
%Book Chapter
\nocite{Tra12}
\nocite{mTX09}
%National Journal
\nocite{FRT16b}
\nocite{HKC13}
\nocite{CAC12}
%National Conference
\nocite{HDT13}
\nocite{CDRa12}
\nocite{HKCa12}
\nocite{CAT11}
\nocite{HTVa11}
\nocite{CATV10}
\nocite{ACT10}

\newpage

\section*{Remerciements}
Nous tenons à exprimer notre profonde gratitude à l’enseignante Karine Serrier pour son soutien, ses conseils et sa patience tout au long de la réalisation de ce rapport.Nous sommes reconnaissants de l’avoir eu comme encadrante et nous espérons que sa passion des mathématiques et sa bonne humeur se ressentiront dans le travail fourni.


	
\chapter{Introduction}

Les modèles déterministes et les modèles probabilistes. Les modèles déterministes sont plus faciles à étudier, mais ils ne reflètent pas toujours la réalité. Les modèles probabilistes, quant à eux, sont une façon naturelle de modéliser l’évolution d’une épidémie, car chaque individu a une certaine probabilité d’être infecté par la maladie.

L’étude de ces problèmes stochastiques est importante pour déterminer si, quand la taille de la population augmente, ils convergent vers un problème déterministe. Nous nous intéresserons plus particulièrement aux épidémies qui se transmettent d’individu en individu et confèrent une immunité aux contacts suivants en cas de rétablissement. Les maladies infantiles, sexuellement transmissibles ou même la grippe correspondent à ce genre de pathologie. On appelle ce type de modèle une épidémie SIR (Susceptible, Infected, Recovered). La population est aisnsi donc séparée en trois catégories.Nous avons les susceptibles qui sont les individus sains étant sensibles à la maladie, les infectés qui sont atteints par l’épidémie et pouvant donc la transmettre pendant une certaine période de contagion, et ceux qui sont immunisés contre une nouvelle infection après avoir été infectés.
\chapter{Etat de l'art}
\section{ceci est une section}
\subsection{ceci est une sous-section}

Etat de l'art
\chapter{Contribution...}

Ma contribution
\chapter{Conclusion}

Mes conclusions


 \listoffigures
\listoftables      
\renewcommand{\bibname}{Bibliography}
\bibliographystyle{alpha}
\bibliography{biblio}


\newpage
\thispagestyle{empty}
\begin{center}
	\textsc{{\bf Abstract}}
\end{center} 

In English
\begin{center}
	\textsc{{\bf Résumé}}
\end{center}

En Français
\end{document}